\documentclass[11pt,a4paper]{article}
\usepackage[utf8]{inputenc}
\usepackage[T1]{fontenc}
\usepackage[margin=2.2cm]{geometry}
\usepackage{amsmath,amssymb,amsthm}
\usepackage{booktabs,longtable,array}
\usepackage{enumitem}
\usepackage{xcolor}
\usepackage[colorlinks=true,linkcolor=blue!50!black,urlcolor=blue!50!black,citecolor=blue!50!black]{hyperref}
\setlength{\parskip}{3pt}
\newtheorem{theorem}{Théorème}[section]
\newtheorem{proposition}[theorem]{Proposition}
\newtheorem{lemma}[theorem]{Lemme}
\newtheorem{corollary}[theorem]{Corollaire}
\theoremstyle{definition}
\newtheorem{question}[theorem]{Question}
\newtheorem{verif}[theorem]{À vérifier}
\theoremstyle{remark}
\newtheorem{remark}[theorem]{Remarque}
\newcommand{\F}{\mathbb{F}}
\newcommand{\Z}{\mathbb{Z}}
\newcommand{\Q}{\mathbb{Q}}
\newcommand{\R}{\mathbb{R}}
\newcommand{\Fix}{\operatorname{Fix}}
\newcommand{\rg}{\operatorname{rang}}
\newcommand{\tr}{\operatorname{tr}}
\newcommand{\srg}{\operatorname{srg}}
\newcommand{\Aut}{\operatorname{Aut}}
\newcommand{\Spec}{\operatorname{Spec}}
\newcommand{\Sh}{\mathrm{Sh}}
\newcommand{\Sout}{S_{\mathrm{out}}}
\newcommand{\Qout}{Q_{\mathrm{out}}}
\newcommand{\st}[1]{\textbf{[#1]}}
\newcommand{\code}[1]{\texttt{#1}}

\title{Le graphe de Conway $\srg(99,14,1,2)$, branche « involution »\\
\large Dossier pour les chercheurs : ce qui est établi, ce qui reste à démontrer, ce qui nous aiderait}
\author{Rédigé par Claude (Fable 5.1) pour D.~Djematene\\
\small à partir d'une lecture intégrale du dépôt \code{amazing-source/conway99-involution} (commit \code{338c00e}, 22 septembre 2026)\\
\small Révisé le 22 septembre 2026 au soir après un second audit indépendant (Claude Opus 5.5) : voir \code{notes/AUDIT\_2026-09-22.md}}
\date{22 septembre 2026 (révision 2)}

\begin{document}
\maketitle

\begin{abstract}
Ce document s'adresse à des chercheurs en combinatoire algébrique et en calcul formel. Il fait le point, sans rien omettre, sur le dépôt de recherche
consacré à la branche « involution » du problème d'existence du graphe de Conway. Il distingue soigneusement (i) ce qui est cité à une source primaire,
(ii) ce qui est démontré dans le dépôt et relu ici, (iii) ce qui est démontré dans ce document même (et donc à relire), (iv) ce qui a été calculé,
et (v) ce qui est douteux, non reproductible ou simplement faux dans la rédaction du dépôt. Les deux dernières parties listent, par ordre de priorité,
ce qu'un chercheur devrait \emph{vérifier} et ce que nous aimerions qu'il \emph{démontre} pour faire avancer la branche.
\textbf{Rien, ni dans le dépôt ni ici, ne tranche l'existence de $\srg(99,14,1,2)$.}
\end{abstract}

\begin{center}
\fbox{\parbox{0.94\textwidth}{\textbf{AVERTISSEMENT --- DOCUMENT EXTERNE AU DÉPÔT.}
Ce dossier a été rédigé le 22 septembre 2026 par une session \emph{différente} de celle qui a produit le dépôt, sur une autre machine, sans accès aux CNF
ni aux sorties intermédiaires de la session d'origine. Il est un \emph{audit}, pas une partie du travail original.
Avant de tenir pour acquis quoi que ce soit d'ici, toute session ou toute personne qui poursuit le travail doit :
(1) ne considérer comme établis que les énoncés étiquetés \st{S} ou \st{R}, qui sont ceux du dépôt, relus ici ;
(2) traiter les énoncés étiquetés \st{N} (spectre de $\Qout$, analogues mod 3, code LCD, remarque de statut du théorème $\F_2$)
comme des énoncés à relire avant tout usage. \emph{Révision 2 :} un second audit a relu chacun d'eux ; le spectre de $\Qout$, le rang $\F_2$ et le rang $\F_3$ (i)
sont corrects et leurs analogues à $k=22$ sont vérifiés numériquement sur un vrai graphe ($\srg(243,22,1,2)$, \code{code/z2/control\_bvls.py}) ;
l'ancienne « parité des $C_4$ » était vide et a été retirée (remarque~\ref{rem:c4}) ; la partie « $Q-3I$ sur $\Z_3$ » de la proposition~\ref{prop:F3}(iv) n'était pas démontrée et a été retirée ;
(3) lire la section~\ref{sec:anomalies} et exécuter le test de \code{notes/ONE10\_PROVENANCE.md} avant de citer un quelconque résultat de solveur ;
(4) signaler dans le dépôt tout désaccord avec ce dossier plutôt que de le corriger silencieusement.
Rien ici ne tranche l'existence de $\srg(99,14,1,2)$ et aucun cas de la branche involution n'est éliminé.}}
\end{center}

\tableofcontents

\section*{Conventions de statut}
Chaque énoncé porte l'une des étiquettes suivantes.
\begin{description}[nosep]
\item[\st{S}] cité à une source primaire, obtenue et lue pendant la session à l'origine du dépôt (les références sont dans \code{notes/SOURCES.md}).
\item[\st{S?}] cité à une source \emph{non obtenue} (résumé de seconde main, ou thèse à la place de l'article).
\item[\st{R}] démontré dans le dépôt ; j'ai relu et refait la preuve ici, elle tient.
\item[\st{N}] démontré \emph{dans ce document}, absent du dépôt ; à relire par un humain avant tout usage.
\item[\st{C}] résultat de calcul du dépôt ; je précise s'il est reproductible avec les fichiers présents.
\item[\st{T}] note tierce (produite par ChatGPT, transmise par l'opérateur) vérifiée dans le dépôt par un script indépendant.
\item[\st{?}] point douteux, incohérence ou lacune que j'ai relevée.
\end{description}

\section{Résumé pour lecteur pressé}
\begin{enumerate}[nosep]
\item Soit $\Gamma$ un hypothétique $\srg(99,14,1,2)$. La littérature (§\ref{sec:biblio}) donne $|\Aut(\Gamma)|\in\{1,2,3\}$. Le cas d'ordre 3 (sans point fixe) est
attaqué par plusieurs projets publics ; \textbf{le cas d'ordre 2 n'a, à notre connaissance, jamais été calculé} (la table de Behbahani commence à $p=3$).
\item Le dépôt montre \st{R} que ce cas est entièrement rigide : l'involution $t$ fixe un seul sommet $x$, agit sur $N(x)=7K_2$ en retournant chaque paire,
aucun sommet extérieur n'est adjacent à son image, et les $\mu$-graphes des paires $\{u,tu\}$ forment un couplage parfait $D$ sur les 42 orbites extérieures.
Il en résulte un modèle exact à $1722$ bits (§\ref{sec:modele}), validé sur $\srg(9,4,1,2)$ (SAT, relevé et vérifié) et $\srg(19,6,1,2)$ (UNSAT immédiat).
\item Le dépôt démontre \st{R} une condition de rang sur $\F_2$ (Théorème~\ref{thm:F2}). Je montre ici \st{N} qu'elle est une \emph{conséquence} de l'identité de quotient
(le dépôt la présente comme « information supplémentaire », ce qui est inexact), j'en donne l'analogue mod 3 (Proposition~\ref{prop:F3}),
et je détermine exactement le spectre du bloc extérieur de la matrice d'orbites (Proposition~\ref{prop:specQout}). Ce sont des invariants utiles pour un générateur,
pas des obstructions nouvelles. \emph{(Révision 2 : une « congruence nouvelle » sur le nombre de $C_4$ induits figurait ici ; elle était vide, voir la remarque~\ref{rem:c4}.)}
\item Tous les calculs exacts (CaDiCaL, Kissat, CP-SAT ; jusqu'à $8$~h par instance) sont \textbf{indécis} ; les heuristiques (recuit) plafonnent loin de zéro.
Un dépassement de temps n'est pas une preuve.
\item J'ai relevé \st{?} plusieurs points à corriger dans le dépôt, dont un log \code{s UNSATISFIABLE} non documenté et non reproductible (§\ref{sec:anomalies}) :
il ne doit \emph{pas} être cité comme un résultat.
\item \emph{(Révision 2)} Le modèle exact est désormais validé sur un cas non trivial : le graphe de Berlekamp--van Lint--Seidel $\srg(243,22,1,2)$ possède une involution
à un seul point fixe ($x\mapsto-x$) et satisfait (R),(D),(I),(P) à $m=11$ sans aucune violation. Les anciens contrôles ($k=4$ : 3 variables ; $k=6$ : non-existence connue)
n'éprouvaient presque rien.
\item Ce dont nous avons besoin (§\ref{sec:verif}--\ref{sec:demontrer}) : une relecture humaine des lemmes (une heure), une réponse à la question
« quelqu'un a-t-il déjà traité $p=2$ ? », et surtout l'énumération des matrices d'orbites $\Qout$ (question~\ref{q:orbit}) puis un \emph{générateur ordonné}
sur les 42 orbites avec entrelacement spectral, ou toute preuve d'exclusion d'un des trois cas du découpage.
\end{enumerate}

\section{Cadre et notations}\label{sec:notations}
$\Gamma$ désigne un $\srg(99,14,1,2)$ hypothétique, $A$ sa matrice d'adjacence :
\[A^2+A-12I=2J,\qquad \Spec(A)=\{14^1,\,3^{54},\,(-4)^{44}\}.\]
Les projecteurs spectraux sont
\[E_3=\tfrac{4I+A}{7}-\tfrac{2J}{77},\qquad E_{-4}=\tfrac{3I-A}{7}+\tfrac{J}{63},\qquad E_{14}=\tfrac{J}{99}\]
(vérifiés : $(A-14I)(A+4I)=-44I-11A+2J$, divisé par $-77$). $E_3$ est $p$-entier pour $p\notin\{7,11\}$ ; $E_{-4}$ pour $p\notin\{3,7\}$.

\paragraph{Réduction enracinée (classique ; écrite notamment par Cesarz--Woldar 2023, Keramatipour 2026, Thakkar--Severini 2026).} On fixe un sommet $x$. Comme $\lambda=1$, $N(x)$ induit $7K_2$, de blocs $B_i=\{(i,0),(i,1)\}$, $i=1,\dots,7$.
Chaque sommet à distance 2 a exactement $\mu=2$ voisins dans $N(x)$, dans deux blocs distincts ; deux sommets extérieurs distincts n'ont pas la même paire
(sinon cette paire non adjacente aurait trois voisins communs). Les 84 sommets extérieurs sont donc en bijection avec les 84 \emph{étiquettes}
$\{(i,s),(j,e)\}$, $i<j$. Toutes les adjacences x--intérieur, intérieur--intérieur et intérieur--extérieur sont forcées ; l'inconnue est le graphe $H$ des adjacences
extérieur--extérieur, $12$-régulier.

\paragraph{Involution.} $t\in\Aut(\Gamma)$, $t^2=1$. On note $\sigma$ le retournement $(i,s)\mapsto(i,1-s)$ de tous les blocs, étendu aux étiquettes par
$\sigma\{a,b\}=\{a',b'\}$ ($a'$ = partenaire de $a$).

\paragraph{Orbites.} Les 42 orbites extérieures sont $U=(i,j,e)$, $i<j$, $e\in\{0,1\}$, de représentant $u=\{(i,0),(j,e)\}$ et $\sigma u=\{(i,1),(j,1-e)\}$.
Pour $U\neq V$ on pose $(a_{UV},b_{UV})=(A[u,v],A[u,\sigma v])\in\{0,1\}^2$ ; ces deux bits sont symétriques en $(U,V)$.
$\Qout=(a+b)$ est le bloc extérieur de la \emph{matrice d'orbites} ordinaire $Q$ ($50\times50$), $\Sout=(a-b)$ celui de la \emph{matrice signée} $S$ ($49\times49$).
$\Gamma_1$ est le graphe des adjacences simples ($a+b=1$), $D$ le couplage des adjacences doubles ($a=b=1$) ; $\Qout=\Gamma_1+2D$.
$N_b\in\{0,1\}^{42\times7}$ est l'incidence orbite/bloc, $\Sh=N_bN_b^{\mathsf T}$ (nombre de blocs partagés, $\Sh_{UU}=2$),
$C\in\{0,\pm1\}^{42\times7}$ la version signée ($C_{U,i}=1$, $C_{U,j}=(-1)^e$), $G=CC^{\mathsf T}$.
On a $C^{\mathsf T}C=12I_7$, $N_b^{\mathsf T}N_b\equiv0\pmod2$, $\rg_{\F_2}N_b=6$ (vérifié numériquement).

\paragraph{Structure de $T(7)$.} Les 42 orbites se projettent 2 pour 1 sur les 21 arêtes de $K_7$ (les paires de blocs) ; deux orbites de même paire sont dites
\emph{sœurs}. $\Sh=J_2\otimes(2I_{21}+A_{T(7)})$ où $T(7)$ est le graphe triangulaire, de spectre $\{10^1,3^6,(-2)^{14}\}$ ; donc
$\Spec(\Sh)=\{24^1,10^6,0^{35}\}$ et $\Spec(G)=\{12^7,0^{35}\}$ (vérifiés numériquement).

\section{État bibliographique (tel qu'établi par le dépôt)}\label{sec:biblio}
\begin{longtable}{p{4.6cm}p{7.6cm}p{2.6cm}}
\toprule
Source & Énoncé utilisé & Statut\\
\midrule
Makhnev--Minakova 2004, Diskr.\ Mat.\ 16:1, th.~1 & $g$ d'ordre premier $\Rightarrow$ $\Fix(g)$ = un sommet ($|g|\in\{2,7\}$), vide ($|g|\in\{3,11\}$) ou un triangle ($|g|=3$) ; $|\Aut|$ divise $2\cdot3^3\cdot7\cdot11$. & \st{S} (original russe lu)\\
Makhnev, notes de Bielefeld & Prop.~2 : sept candidats pour $\Fix(t)$ ; seul le cas « un sommet » rend $\chi_2(t)=(4\alpha_0+\alpha_1-18)/7$ entier. Lemme (même page) : $|\Fix(t)|=1+|\Fix(t)\cap[a]|+2l$, $l$ = nombre de $t$-orbites de taille 2 sur les arêtes de $[a]$, $a\in\Fix(t)$ ; avec $\Fix(t)=\{x\}$ cela donne $l=0$, c'est-à-dire le lemme~\ref{lem:flip}. & \st{S}\\
Behbahani 2009 (thèse, th.~4.14, table 21) & Aucune matrice d'orbites pour $p\in\{5,7,11\}$ (tous nombres de points fixes admissibles) ni pour $p=3$ avec points fixes. \textbf{$p=2$ absent de la table.} & \st{S} (thèse) / \st{S?} (article 2011 non obtenu)\\
Crnković--Maksimović 2020 & Résumé publié : « we rule out automorphism groups isomorphic to $Z_6$, $S_3$, $Z_9$, or $E_9$ » ; le résumé ne mentionne pas de matrices d'orbites pour une involution seule. & \st{S?} (résumé lu, article non obtenu)\\
Keramatipour 2026, arXiv:2604.23037 & Réduction enracinée + solveurs SAT sur $\srg(99,14,1,2)$ sans symétrie imposée : le résumé conclut à « the incapability of SAT solvers facing this problem in a reasonable time ». & \st{S} (résumé)\\
Cesarz--Woldar 2025 & Sans ordinateur : $2\mid|G|\Rightarrow|G|\mid6$ ; $7\mid|G|\Rightarrow G\cong\Z_7$. & \st{S}\\
Ishida 2026, arXiv:2606.29183 & Identité trace--rang ; redonne $a_1(t)=14$ pour une involution ; pas d'exclusion. Prop.~8.7 : $a_1(g)=18$ pour $g$ d'ordre 3 sans point fixe. & \st{S}\\
Thakkar--Severini 2026, arXiv:2608.11211v2 & Réduction enracinée ; écrivent que le cas $\Z_7$ à un point fixe « reste ouvert ». & \st{S}\\
Reimbayev 2026, arXiv:2608.19410v1 & Recensement des 208 sous-graphes induits d'ordre 7 dans $\srg(n,k,1,2)$. & \st{S} (source \TeX{} auditée)\\
Dépôt public \code{YesterdaysLemon/conway-99-research} & $708\le n_3\le4158$ (vague 23) ; famille affine d'ordre 7 ; $\rg_{\F_7}(M)\ge27$ pour $M=B^{\mathsf T}(3I-A)B+J$. & \st{S} (dépôt lu)\\
\bottomrule
\end{longtable}
Conséquence : $|\Aut(\Gamma)|\in\{1,2,3\}$, et une involution engendre tout le groupe.
Deux corrections de lecture faites par le dépôt : le cas $\Z_7$ à un point fixe est clos depuis 2009/2011 (Thakkar--Severini le disent ouvert) ; les vagues 69/72 du dépôt public (ordre 11) sont redondantes avec Behbahani--Lam.

\begin{verif}\label{ver:biblio}
\emph{Révision 2 :} « Crnković--Maksimović excluent les deux groupes d'ordre 6 » est confirmé par leur résumé ($Z_6$ et $S_3$). Reste à confirmer, sur les textes
complets : (a) que Behbahani--Lam 2011 ne traite pas $p=2$ ; (b) que l'article de 2020 ne contient pas d'énumération de matrices d'orbites pour l'involution seule
(dans l'analyse de $\Z_6$ ou de $S_3$). Si (b) est faux, la branche est peut-être déjà partiellement traitée.
\end{verif}

\section{Ce qui est démontré}\label{sec:demontre}

\subsection{Rigidité de l'involution \st{R}}
\noindent\emph{Antériorité (révision 2).} Le lemme~\ref{lem:flip} découle directement du lemme de Makhnev cité au §\ref{sec:biblio} ($l=0$) ; le lemme~\ref{lem:noself}
reprend $a_1(t)=14$ de Makhnev--Minakova. Ce qui semble propre au dépôt est la proposition~\ref{prop:D} et le modèle exact qui en découle ; ne pas présenter
la « rigidité » dans son ensemble comme nouvelle.
\begin{lemma}[l'involution est le retournement]\label{lem:flip}
$t=\sigma$ sur $N(x)$, et $t$ agit sur les étiquettes par $\sigma$.
\end{lemma}
\begin{proof}
$\Fix(t)=\{x\}$ par Makhnev--Minakova. $t$ préserve $N(x)$ et son couplage, donc permute les blocs ; $t^2=1$, donc chaque bloc est fixé ou échangé avec un autre.
Un bloc fixé est retourné (aucun sommet intérieur n'est fixe). Si $t$ échangeait $B_i$ et $B_j$, on aurait $t(a)=b$, $t(b)=a$ avec $a\in B_i$, $b\in B_j$, et l'étiquette
$\{a,b\}$, qui est un sommet extérieur, serait fixe : contradiction.
\end{proof}

\begin{lemma}[aucun extérieur n'est adjacent à son image]\label{lem:noself}
Si $u\sim tu$ alors $u\in N(x)$. Donc $a_1(t)=14$, et $\chi_3(t)=\chi_{-4}(t)=0$.
\end{lemma}
\begin{proof}
L'arête $\{u,tu\}$ est $t$-invariante ; son unique voisin commun ($\lambda=1$) est donc $t$-fixe, c'est $x$. Avec $a_0=1$, $a_1=14$, $a_2=84$ :
$\chi_3(t)=\tr(tE_3)=\tfrac{42}{77}+14\cdot\tfrac{9}{77}-84\cdot\tfrac{2}{77}=0$ et $\chi_{-4}(t)=\tr(t)-\chi_{14}(t)-\chi_3(t)=1-1-0=0$.
\end{proof}

\begin{proposition}[le couplage $D$]\label{prop:D}
Pour tout extérieur $u$, $N(u)\cap N(\sigma u)$ est une orbite $\{v,\sigma v\}=:D(u)$ d'extérieurs, et $D$ est un couplage parfait sur les 42 orbites.
Chaque extérieur est adjacent aux deux membres d'exactement une orbite.
\end{proposition}
\begin{proof}
$u\not\sim\sigma u$, donc $|N(u)\cap N(\sigma u)|=\mu=2$ ; l'ensemble est $\sigma$-stable sans point fixe, c'est une orbite ; elle est extérieure car les voisins
intérieurs de $u$ et de $\sigma u$ sont disjoints et $x\not\sim u$. Si $D(u)=\{v,\sigma v\}$ alors $u,\sigma u\in N(v)\cap N(\sigma v)$, d'où $D(v)=\{u,\sigma u\}$.
\end{proof}

\subsection{Le modèle exact à 42 orbites \st{R}}\label{sec:modele}
\begin{theorem}[équivalence exacte]\label{thm:modele}
Étant donné la structure forcée, une matrice symétrique $0/1$ à diagonale nulle, $\sigma$-invariante, avec $A[u,\sigma u]=0$ pour $u$ extérieur, vérifie $A^2+A-12I=2J$
si et seulement si les bits $(a_{UV},b_{UV})$ satisfont :
\begin{itemize}[nosep]
\item[(R)] $\sum_{V\ne U}(a_{UV}+b_{UV})=12$ ;
\item[(D)] il existe exactement un $V$ avec $a_{UV}=b_{UV}=1$ ;
\item[(I)] pour tout intérieur $c$ : $\#\{\text{voisins } w \text{ de } u \text{ dans } H : c\in w\}=1$ si $c\in u$ ou $c'\in u$, et $2$ sinon ;
\item[(P)] pour $U\ne V$ : $\sum_W(a_{UW}a_{VW}+b_{UW}b_{VW})+a_{UV}=2-|u\cap v|$ et $\sum_W(a_{UW}b_{VW}+b_{UW}a_{VW})+b_{UV}=2-|u\cap\sigma v|$.
\end{itemize}
\end{theorem}
\begin{proof}[Preuve (détaillée ici, esquissée dans le dépôt)]
L'identité $A^2+A-12I=2J$ équivaut à : degré 14 partout, $\lambda=1$, $\mu=2$. Les paires de sommets sont de sept types.
(1) $\{x,c\}$, $c$ intérieur : voisin commun unique $c'$. (2) $\{x,u\}$ : voisins communs = les deux intérieurs de $u$. (3) $\{c,c'\}$ : voisin commun $x$.
(4) $\{c,d\}$ de blocs différents : voisins communs = les étiquettes $\{c,d\}$... précisément l'unique étiquette $\{c,d\}$ et $x$, soit 2. Ces quatre types sont forcés.
(5) $\{c,u\}$ intérieur--extérieur : les voisins communs intérieurs sont $c'$ si $c'\in u$ (et jamais $c$ lui-même) ; $x$ n'en est jamais un ; les voisins communs extérieurs
sont les $w\sim u$ avec $c\in w$. Si $c\in u$ ($c\sim u$) il en faut $1$ au total, et $c'\notin u$, donc exactement $1$ extérieur. Si $c'\in u$ ($c\not\sim u$) il en faut $2$, dont $c'$ : exactement $1$ extérieur.
Sinon il en faut $2$, tous extérieurs. C'est (I), en comptant $w=v$ via $a_{UV}$ ($c\in v$) et $w=\sigma v$ via $b_{UV}$ ($c\in\sigma v$).
(6) $\{u,\sigma u\}$ : non adjacents, voisins communs intérieurs : aucun ; extérieurs : pour chaque orbite $W$, $w$ compte si $a_{UW}b_{UW}=1$ et $\sigma w$ aussi ; total $2\sum_W a_{UW}b_{UW}=2$. C'est (D).
(7) $\{u,v\}$, $V\ne U$ : voisins communs intérieurs $|u\cap v|$ ; extérieurs : $\sum_{W\ne U,V}(a_{UW}a_{VW}+b_{UW}b_{VW})$ (les sommets $\sigma u,\sigma v$ ne contribuent pas car $A[u,\sigma u]=0$ et $A[v,\sigma v]=0$) ; le total doit valoir $2-a_{UV}$. Même calcul pour $\{u,\sigma v\}$ avec les termes croisés. C'est (P).
Le degré : $2+\sum_V(a+b)=14$, c'est (R). Les paires $\{\sigma u,\sigma v\}$ et $\{\sigma u,v\}$ sont les images des précédentes.
\end{proof}
Les hypothèses « $\sigma$-invariante » et « $A[u,\sigma u]=0$ » sont justifiées par les lemmes~\ref{lem:flip} et \ref{lem:noself} ; elles ne sont pas des variables du modèle.
Le modèle est implémenté dans \code{code/z2/orbit\_sat.py} (CNF, encodage \emph{seqcounter} de PySAT) et \code{cpsat\_orbit.py} (OR-Tools), pour $k=2m$ quelconque, avec relevé en $99\times99$ et vérification entière de l'identité.
Contrôles \st{C} : $m=2$ ($\srg(9,4,1,2)$, le graphe de la tour $3\times3$, qui possède bien une involution à un point fixe) : SAT, relevé vérifié ; $m=3$ ($\srg(19,6,1,2)$ n'existe pas, valeurs propres irrationnelles) : UNSAT immédiat.
\emph{Révision 2 :} ces deux contrôles sont presque vides (à $m=2$ le modèle n'a que 3 variables ; un UNSAT à $m=3$ ne teste pas qu'une solution existante est acceptée).
Contrôle fort ajouté : le graphe de Berlekamp--van Lint--Seidel $\srg(243,22,1,2)$ avec l'involution $x\mapsto-x$ (un seul point fixe) est l'instance $m=11$ du même modèle ;
sa table $(a,b)$ satisfait (R),(D),(I),(P) sans violation, ainsi que les équations de $\Qout$ et de $\Sout$ (\code{code/z2/control\_bvls.py}, mode léger, vérifié le 22/09).
Le mode \code{--cnf} (vérification de l'encodage CNF lui-même à $m=11$ : 12,6 M variables, 28,3 M clauses) reste à exécuter sur la machine de calcul.
Pour $m=7$ : 1722 variables primaires, 666\,183 variables CNF, 1\,510\,448 clauses (1\,510\,446 sans découpage) ; regénération vérifiée le 22/09, empreintes SHA-256 dans \code{LAB\_INSTRUCTIONS.md}.

\subsection{Les deux quotients \st{R}}
\begin{proposition}\label{prop:quot}
(i) $Q^2+Q-12I=2\hat J$ avec $\hat J_{VU}=|U|$ ; en particulier $\Qout^2+\Qout=12I+4J-\Sh$ et $N_b^{\mathsf T}\Qout=4J-2N_b^{\mathsf T}$ (sommes par bloc : 2 sur les blocs de $U$, 4 sur les autres).
(ii) Dans la base $f_U=e_u-e_{\sigma u}$ de l'espace anti-invariant (dimension 49), $A$ agit par $S=\begin{pmatrix}-I_7&C^{\mathsf T}\\ C&\Sout\end{pmatrix}$,
avec $S^2+S=12I$ exactement, $\tr S=-7$, donc $\Spec(S)=\{3^{27},(-4)^{22}\}$.
(iii) $C^{\mathsf T}\Sout=0$, $\Sout^2+\Sout=12I-G$, $\Sout\in\{0,\pm1\}^{42\times42}$ avec dix entrées non nulles par ligne, $\Spec(\Sout)=\{0^7,3^{20},(-4)^{15}\}$,
et $P_3:=(\Sout+4I-\tfrac13G)/7$ est le projecteur orthogonal de rang 20 sur l'espace propre $3$.
\end{proposition}
\begin{proof}
(i) est la projection standard de l'identité SRG sur les orbites ; les sommes par bloc sont le bloc (intérieur, extérieur) de $Q^2+Q$.
(ii) $J$ tue les vecteurs anti-invariants (leurs coordonnées se compensent sur chaque orbite), d'où $S^2+S=12I$ ; les entrées diagonales sont $-A[u,\sigma u]$ : $-1$ pour les sept orbites intérieures, $0$ ailleurs (lemme~\ref{lem:noself}) ; $3f-4g=-7$, $f+g=49$ donne $f=27$.
(iii) blocs de $S^2+S=12I$. $\Sout$ tue les colonnes de $C$ ; sur leur orthogonal (dimension 35) $G=0$, donc $\Sout^2+\Sout=12I$ et la trace nulle force $3f'-4g'=0$, $f'+g'=35$, $f'=20$.
\end{proof}

\begin{proposition}[spectre exact de $\Qout$] \st{N}\label{prop:specQout}
$\Qout$, $\Sh$ et $J$ commutent, et $\Spec(\Qout)=\{12^1,(-2)^6,3^{20},(-4)^{15}\}$. Les vecteurs propres pour $-2$ sont les $\mathbf 1_2\otimes v$, $v$ dans l'espace propre $3$ de $T(7)$, c'est-à-dire les fonctions $U=(i,j,e)\mapsto y_i+y_j$ avec $\sum y=0$.
\end{proposition}
\begin{proof}
De $N_b^{\mathsf T}\Qout=4J-2N_b^{\mathsf T}$ on tire $\Qout\Sh=8J-2\Sh=\Sh\Qout$. Sur chaque espace propre commun de $(\Sh,J)$, $q^2+q=12+4j-s$ :
sur $\mathbf 1$ ($s=24$, $j=42$) : $q=12$ ; sur l'espace $s=10$ (dimension 6, $j=0$) : $q^2+q=2$, $q\in\{1,-2\}$ ; sur $s=0$ (dimension 35) : $q\in\{3,-4\}$.
Avec $a$ valeurs $1$ et $b$ valeurs $3$, $\tr\Qout=0$ donne $12+a-2(6-a)+3b-4(35-b)=0$, soit $3a+7b=140$ ; les seules solutions avec $0\le a\le6$, $0\le b\le35$ sont $(a,b)=(0,20)$.
Enfin la condition « $(\Qout+2I)(\mathbf 1_2\otimes v)=0$ » se réduit, pour $v_e=y_i+y_j$, aux sommes par bloc de (i) : elle n'ajoute rien au modèle, mais le spectre entier est utilisable pour l'entrelacement.
\end{proof}
\noindent\emph{Contrôle (révision 2).} La même preuve donne, pour $k=2m$ quelconque, $\Spec(\Qout)=\{k-2,\,(-2)^{m-1},\,r^{b},\,s^{c}\}$ avec $b,c$ fixés par $\tr\Qout=0$.
Sur $\srg(243,22,1,2)$ ($m=11$, $r=4$, $s=-5$) la prédiction $\{20,(-2)^{10},4^{55},(-5)^{44}\}$ est exactement le spectre calculé ; de même l'analogue du
théorème~\ref{thm:F2} ($\rg_{\F_2}\bar S=55$, $\rg_{\F_2}(\bar\Gamma+N_bN_b^{\mathsf T})=44$) est vérifié (\code{control\_bvls.py}).

\subsection{La condition de rang sur $\F_2$ \st{R}, et son statut logique \st{N}}
\begin{theorem}[dépôt, rapport §3]\label{thm:F2}
Soit $\bar\Gamma=\Gamma_1\bmod2\in\F_2^{42\times42}$. Alors $\bar\Gamma^2+\bar\Gamma=N_bN_b^{\mathsf T}$, $\bar\Gamma N_b=0$, et
$K:=\bar\Gamma^2=\bar\Gamma+N_bN_b^{\mathsf T}$ est un idempotent symétrique de rang \emph{exactement} $20$ avec $KN_b=0$.
\end{theorem}
\begin{proof}
$S$ est entière de polynôme caractéristique $(x-3)^{27}(x+4)^{22}$ ; modulo 2 il devient $(x+1)^{27}x^{22}$, et $\bar S^2=\bar S$ ; donc $\bar S$ est diagonalisable de rang 27.
Comme $\bar S=\begin{pmatrix}I_7&N_b^{\mathsf T}\\ N_b&\bar\Gamma\end{pmatrix}$ ($\Sout\equiv\Qout\equiv\bar\Gamma$ mod 2), un complément de Schur donne $\rg(\bar\Gamma+N_bN_b^{\mathsf T})=20$.
Les identités polynomiales sont les réductions mod 2 des blocs de $\bar S^2=\bar S$.
\end{proof}
\begin{remark}[statut logique]\label{rem:statut}
Le dépôt écrit que la valeur 20 « n'est pas impliquée par les identités de comptage » et est « une information véritablement supplémentaire ». C'est vrai des seules identités mod 2 ($K^2=K$ ne donne que $\rg K$ pair), mais la preuve ci-dessus montre que le théorème est une \emph{conséquence} de l'identité entière $S^2+S=12I$ et de $\tr S=-7$, elles-mêmes conséquences de (R),(D),(I),(P). Il ne s'agit donc pas d'une obstruction indépendante du modèle, mais d'un \emph{invariant global} : son intérêt est (a) de s'appliquer à toute relaxation qui ne contiendrait que le quotient signé, et surtout (b) de fournir un filtre sur les matrices \emph{partielles} : toute sous-matrice principale de $\bar\Gamma+N_bN_b^{\mathsf T}$ a un $\F_2$-rang $\le20$. La note \code{NOTE\_involution\_case.md} et le \code{README} gardent encore la formulation initiale « structure 2-adique $\Rightarrow$ » ; à harmoniser avec le rapport.
\end{remark}
\begin{remark}[reformulation en code] \st{N}
$K$ étant un idempotent symétrique, $L:=\operatorname{Im}K$ est un sous-espace de dimension 20 de $\F_2^{42}$ tel que $\F_2^{42}=L\oplus L^\perp$ (un code \emph{LCD} $[42,20]$), et $KN_b=0$ signifie $L\subseteq(\operatorname{col}N_b)^\perp$, sous-espace de dimension 36. Réciproquement $\bar\Gamma=K+N_bN_b^{\mathsf T}$ : le graphe des adjacences simples, réduit mod 2, est entièrement déterminé par le code $L$. Question ouverte : les contraintes combinatoires ($\bar\Gamma$ est une matrice $0/1$ 10-régulière à diagonale nulle) se traduisent-elles en contraintes sur l'énumérateur de poids de $L$ (identités de MacWilliams) ?
\end{remark}

\subsection{Analogues $p$-locaux \st{N}}
\begin{proposition}\label{prop:F3}
(i) $\rg_{\F_3}(S\bmod3)=22$, de façon équivalente $\rg_{\F_3}(\Sout+G)=15$.
(ii) Dans la base $\{e_x,\;e_u+e_{\sigma u}\}$ de l'espace invariant, $M:=I+Q-\hat J$ réduite mod 3 est un idempotent de rang exactement 27 ; explicitement, en blocs $(x,\text{intérieur},\text{extérieur})$,
\[M\equiv\begin{pmatrix}0&0&\mathbf 1^{\mathsf T}\\ 0&J-I_7&N_b^{\mathsf T}+J\\ -\mathbf 1&N_b+J&I+\Gamma_1-D+J\end{pmatrix}\pmod 3 .\]
(iii) Sur $\F_2$, la condition analogue sur $Q$ est automatiquement impliquée par le théorème~\ref{thm:F2}.
(iv) Sur $\Z_2$ et $\Z_3$, les formes normales de Smith de $S-3I$ et $S+4I$ sont triviales (unités et zéros) ; pour $S$, le contenu $p$-adique ($p\in\{2,3\}$) est donc épuisé par (i) et le théorème~\ref{thm:F2}. Sur $\Z_7$ le contenu n'est pas trivial (§\ref{sec:demontrer}, Q\ref{q:7}).
\emph{(Révision 2 : la version précédente affirmait aussi la trivialité pour $Q-3I$ ; la preuve ne couvre que $S$. Sur $\Z_2$ elle s'étend à $Q$ car $E_{14}=J/99$, $E_3$, $E_{-4}$ sont 2-entiers ; sur $\Z_3$, $E_{14}$ et $E_{-4}$ ne sont pas 3-entiers et l'énoncé pour $Q-3I$ reste non démontré.)}
\end{proposition}
\begin{proof}
(i) $S^2+S=12I\equiv0$ mod 3 : $S$ est diagonalisable sur $\F_3$ de polynôme caractéristique $x^{27}(x+1)^{22}$, donc de rang 22 ; complément de Schur par le bloc $-I_7$ : $\rg(\Sout-C(-I)^{-1}C^{\mathsf T})=\rg(\Sout+G)=15$.
(ii) $E_3$ est 3-entier, commute avec $t$, donc préserve le réseau invariant $\Lambda_e=\Z_3^{99}\cap V^{\mathrm{inv}}$, dont $\{e_x,e_u+e_{\sigma u}\}$ est une base ; $E_3|_{\Lambda_e}$ est un idempotent de rang $(54+\chi_3(t))/2=27$, et sa réduction a le même rang. Dans cette base $A$ agit par $Q$ et $J$ par $\hat J$ ; mod 3, $7^{-1}\equiv1$, $4\equiv1$, $2/77\equiv1$, d'où $M$. Les blocs viennent de $Q_{x,A_i}=2$, $Q_{A_i,x}=Q_{A_i,A_i}=1$, $Q_{A_i,U}=Q_{U,A_i}=[i\in U]$, $\hat J$ de colonnes $|U|$.
(iii) Mod 2, $E_3|_{\Lambda_e}\equiv Q$ dont le rang vaut celui de $[c\,|\,\bar S]$ avec $c=(\mathbf 1_7;0)$ ; or $\bar S c=c$ car chaque orbite rencontre deux blocs, donc le rang est $\rg\bar S=27$.
(iv) Pour $p\in\{2,3\}$ et $\theta\in\{3,-4\}$, $S-\theta I$ est une unité fois un projecteur $p$-entier de rang connu, dont l'image et le noyau sont des facteurs directs.
\end{proof}
\begin{remark}
Comme au théorème~\ref{thm:F2}, ces énoncés sont des conséquences des identités de quotient (voir remarque~\ref{rem:statut}). Leur valeur : (i) porte sur la matrice \emph{signée}, information perdue mod 2 ; (ii) fait intervenir le couplage $D$. Tous deux donnent des filtres de sous-matrices principales ($\F_3$-rang $\le15$, resp.\ $\le27$) ; \emph{attention (révision 2)} : un filtre « rang $\le r$ » ne peut
rejeter aucune sous-matrice de taille $\le r$, il n'agit donc qu'à partir de la taille $16$ (resp.\ $28$ ; $21$ pour le filtre $\F_2$ du théorème~\ref{thm:F2}) dans un générateur ligne par ligne.
Le filtre qui agit tôt est l'entrelacement : toute sous-matrice principale de $\Qout$ a au plus une valeur propre $>3$.
\end{remark}

\subsection{Comptage d'orbites : ce qui reste après correction \st{N}}
\begin{remark}[révision 2 : la « parité des $C_4$ » était vide]\label{rem:c4}
La version précédente énonçait comme résultat nouveau que, dans la branche involution, le nombre $c_4$ de $C_4$ induits est impair et le nombre de diamants $\equiv0\pmod4$.
C'est vrai mais \emph{sans contenu}, pour \emph{tout} $\srg(99,14,1,2)$ : deux sommets non adjacents ont exactement deux voisins communs, et ceux-ci ne sont jamais adjacents
(sinon leur arête serait dans deux triangles, contre $\lambda=1$). Il n'y a donc aucun diamant, chaque non-arête est la diagonale d'exactement un $C_4$ induit, et
$c_4=4158/2=2079$ est une constante (impaire, divisible par 3). Plus généralement, tous les comptes de sous-graphes induits d'ordre 4 sont constants dans un
$\srg(n,k,1,2)$ (le seul compte non déterminé par les paramètres est celui des $K_4$, nul ici) ; le paramètre libre $n_3$ des recensements n'est donc \emph{pas}
un compte d'ordre 4.
Ce qui reste exact et utile : les $C_4$ induits $\sigma$-invariants sont exactement les 21 cycles $u\,v\,\sigma u\,\sigma v$ avec $D(u)=\{v,\sigma v\}$ (un par arête de $D$) ;
la preuve précédente (un $C_4$ invariant ne contient pas $x$, ne peut contenir d'orbite intérieure, et est de cette forme) reste valable.
\end{remark}

\subsection{Découpage en cas \st{R}, vérifié numériquement \st{N}}
Le groupe des relabellisations de la structure forcée est $S_2\wr S_7$ (ordre $645\,120$) ; $\sigma$ en est central, il agit donc sur les tables d'orbites.
Le stabilisateur de l'orbite de référence $U_0=(1,2,0)$ est d'ordre $15\,360$ et a exactement trois orbites sur les 41 autres orbites : $\{(1,2,1)\}$ (sœur, taille 1), celle de $(1,3,0)$ (un bloc commun, taille 20), celle de $(3,4,0)$ (blocs disjoints, taille 20). Le découpage sib/one/dis sur $D(U_0)$ est donc exhaustif et sans redondance (calcul refait ici par énumération des 645\,120 éléments).
Dans le cas dis, le stabilisateur de $(U_0,D(U_0))$ est d'ordre 768 et contient l'échange des blocs 1 et 2, qui envoie « $u_0\sim\{(1,0),(2,1)\}$ » sur « $u_0\sim\{(1,1),(2,0)\}$ » : le sous-découpage dis00/dis10 (non adjacent / adjacent à la sœur) est sain. Dans le cas one, la contrainte (I) sature les deux places du bloc 1 par $D(U_0)$ : $u_0$ n'est adjacent à \emph{aucun} sommet de l'orbite sœur, et le sous-cas « one10 » est vide par propagation unitaire (voir §\ref{sec:anomalies}).

\section{Inventaire exhaustif du dépôt}\label{sec:inventaire}
Dix commits, tous du 21--22 septembre 2026 (19:00:45 $\to$ 02:59:22). Les fichiers \code{*.cnf} sont ignorés par git (regénérables en 2 s, 29 Mo chacun).

\begin{longtable}{p{5.2cm}p{9.6cm}}
\toprule
Fichier & Rôle / état\\
\midrule
\multicolumn{2}{l}{\emph{Racine et notes}}\\
\code{README.md} & Résumé, disposition, commandes de reproduction. \emph{Réécrit en révision 2} (périmètre, corrections, point de départ).\\
\code{LAB\_INSTRUCTIONS.md} & \emph{(Révision 2)} Instructions pour la session de la machine de calcul : contrôles, tests, prochaine tâche.\\
\code{notes/AUDIT\_2026-09-22.md} & \emph{(Révision 2)} Second audit, fichier par fichier.\\
\code{notes/NOTE\_involution\_case.md} & Note de travail : bibliographie, rigidité, modèle, quotients, condition $\F_2$ (avec l'ancienne preuve 2-adique et un préambule donnant la preuve courte), calculs, plan. Cite des chemins \code{scratchpad/z2\_signed/} qui ne sont pas dans le dépôt.\\
\code{notes/RESEARCH\_MEMORY.md} & Mémoire de session (horodatée) ; cohérente avec le reste ; mentionne des sorties non archivées (\code{result\_cpsat\_none.json}, \code{rowfix\_seed*.json}, \code{rows\_U0.npy}).\\
\code{notes/WHAT\_WOULD\_HELP.md} & Liste des besoins (documents, outils, calcul, vérification humaine, questions). Reprise et étendue en §\ref{sec:besoins}.\\
\code{notes/SOURCES.md} & Liens vers les sources primaires ; aucun PDF redistribué.\\
\code{notes/AUDIT\_reimbayev\_order7.md} & Audit du recensement d'ordre 7 (§\ref{sec:audit}).\\
\code{notes/third\_party/*} & \emph{Supprimé en révision 2} (branche d'ordre 3, hors sujet ; intègre, sommes SHA-256 vérifiées ; récupérable au commit \code{9eb634d}).\\
\multicolumn{2}{l}{\emph{Rapports}}\\
\code{report/conway99\_involution\_report.tex/.pdf} & Rapport technique anglais (5 sections + bibliographie) ; preuve courte du théorème $\F_2$. Révision 2 : « new necessary condition / genuinely additional information » et « batch running » corrigés.\\
\code{report/involution\_notes\_fr.tex/.pdf} & Notes françaises pédagogiques ; même contenu mathématique. Révision 2 : mêmes corrections.\\
\code{report/Dockerfile.tex} & Image Ubuntu + TeX Live pour compiler (utilisée pour ce document).\\
\multicolumn{2}{l}{\emph{Code : modèle à orbites (\code{code/z2/})}}\\
\code{orbit\_sat.py} & Modèle CNF exact (R,D,I,P), cas sib/one/dis/none, relevé et vérification. \textbf{Ne connaît pas} les sous-cas dis00/dis10/one00/one10.\\
\code{cpsat\_orbit.py} & Même modèle en CP-SAT ; option \code{--rowfix-seed}.\\
\code{run\_orbit.py} & Lance CaDiCaL (PySAT) avec limite murale, relève tout modèle.\\
\code{row\_fix.py} & Fixe aléatoirement toute la ligne de $U_0$ (ligne valide construite par DFS) et mesure le temps de décision.\\
\code{rows.py} & Énumère les lignes candidates de $U_0$ (motifs $Q$ puis signes) ; le nombre obtenu n'est pas archivé.\\
\code{anneal\_orbit.py} & Recuit simulé « bit-flip » sur les 1722 bits ; énergie = violations de (R),(D),(I),(P).\\
\code{anneal\_struct.py} & Recuit structuré : état = couplage $D$ + graphe 10-régulier $G$ + signes ; (R),(D) vraies par construction.\\
\code{joint\_dfs.py} & DFS ligne par ligne sur la table $(Q,S)$ avec contraintes de Gram exactes ; « bloque aux lignes 2--6 » (mémoire), sortie non archivée.\\
\code{sout\_dfs.py} & DFS sur $\Sout$ seule ((W),(B),(G)) ; sortie non archivée.\\
\code{dfs\_backup.py}, \code{common.py}, \code{cpsat\_model.py} & Relaxation « matrice signée seule », diagonale $0$ ou relâchée ; logs \code{results/k14\_default.log}, \code{k14\_relaxdiag.log}. \st{?} \emph{Révision 2 :} n'impose que le bloc extérieur $\Sout^2+\Sout=12I-G$ (903 équations, cf.\ \code{\#kLinearN: 903} du log) et \emph{pas} $C^{\mathsf T}\Sout=0$ (294 équations) : relaxation plus faible que « $S^2+S=12I$ ». Sans effet sur les résultats (UNKNOWN ; toute solution serait rejetée par \code{verify}).\\
\code{probe99/model.py, run\_cadical.py, run\_cpsat.py, validate\_k4.py} & Modèle à 99 sommets construit génériquement (décomposition automatique des voisins communs), identification $\sigma$ optionnelle, validation exhaustive à $k=4$ (\emph{une} solution étiquetée dans chacun des trois modes, vérifiée, contrôle négatif ; révision 2 : la version précédente disait 3, le « 3 » affiché par le script est le nombre de 4-cycles étiquetés possibles, codé en dur). Défaut mineur : dans \code{run\_cpsat.py}, \code{--no-self-adjacent} est \code{store\_true} avec \code{default=True}, donc non désactivable. Résultats à $k=14$ non archivés (« UNKNOWN » d'après la mémoire).\\
\code{Dockerfile}, \code{run\_batch.sh} & Image \code{satbox} (Kissat 4.0.4 + CaDiCaL compilés) ; script de batch \emph{avec un chemin absolu vers un autre répertoire de travail} (révision 2 : répertoire désormais passé en argument) et ne lançant que 6 des 7 instances résumées.\\
\code{control\_bvls.py} & \emph{(Révision 2)} Contrôle positif fort sur $\srg(243,22,1,2)$ ($m=11$) ; mode léger vérifié, mode \code{--cnf} pour la machine de calcul.\\
\code{anneal\_*.out/.err}, \code{struct\_*.out}, \code{bitflip\_5.out}, \code{*.json} & Traces et bilans des recuits (voir tableau §\ref{sec:calculs}).\\
\multicolumn{2}{l}{\emph{Code : audit d'ordre 7 (\code{code/r7\_check/})}}\\
\code{parse\_r7.py} & Transcription des 208 formules depuis le source \TeX{} (non inclus), identité $\sum z_i=\binom n7$, spécialisation à $(99,14)$, énumération des $(n_3,z_{11})$ admissibles.\\
\code{match705.py}, \code{corrected.py}, \code{family\_n3.py} & Appariement aux classes du dépôt public (non inclus), corrections, reconstruction exacte de la famille affine.\\
\code{spec99.json}, \code{repo\_family\_affine.json}, \code{feasible\_n3\_z11.json} & Données produites (coefficients $(a,b,c)$ des 208 formules ; famille du dépôt ; 4159 valeurs de $n_3$ admissibles).\\
\multicolumn{2}{l}{\emph{Résultats (\code{results/})}}\\
\code{BATCH\_2026-09-21\_overnight.md} & Bilan de la nuit : 7 instances, toutes UNKNOWN après 8 h.\\
\code{batch\_*.log} (10 fichiers) & Sorties brutes (mode \code{-q} : une ligne). Trois d'entre eux ne sont pas dans le bilan (§\ref{sec:anomalies}).\\
\code{result\_cpsat\_one.json} & CP-SAT, cas one, 907 s, UNKNOWN, 2477 conflits.\\
\bottomrule
\end{longtable}

\section{Calculs effectués et résultats}\label{sec:calculs}
\begin{longtable}{p{3.9cm}p{2.6cm}p{2.2cm}p{5.8cm}}
\toprule
Modèle / instance & Solveur & Temps & Résultat \st{C}\\
\midrule
Orbites, $m=2$ & CaDiCaL & $<1$ s & SAT ; relevé $9\times9$ = $\srg(9,4,1,2)$ vérifié.\\
Orbites, $m=3$ & CaDiCaL & $<1$ s & UNSAT (attendu).\\
99 sommets, $k=4$, sans $\sigma$ / avec $\sigma$ / avec $\sigma$ et $u\not\sim\sigma u$ & CaDiCaL (énum.) & s & 1 / 1 / 1 solution, vérifiées ; contrôle négatif réussi (rejoué le 22/09).\\
Orbites, $m=11$ : $\srg(243,22,1,2)$ planté & contrôle direct & $<1$ min & (R),(D),(I),(P) : 0 violation ; équations $\Qout$, $\Sout$, spectre et rang $\F_2$ prédits : exacts (révision 2).\\
Signée seule, $m=7$, diag $0$ & CP-SAT, 6 fils & 1500 s & UNKNOWN ; 936\,300 conflits, 34,7 M branches.\\
Signée seule, diag relâchée & CP-SAT, 8 fils & 900 s & UNKNOWN (log tronqué au début de la recherche).\\
Orbites, ligne de $U_0$ fixée (4 graines) & CaDiCaL & $\ge30$ min CPU & indécis (fichiers \code{rowfix\_seed*.json} non archivés).\\
Orbites, cas one & CP-SAT, 6 fils & 907 s & UNKNOWN ; 2477 conflits, 400\,327 branches.\\
Orbites, sib / one / dis & CaDiCaL natif & 8 h & UNKNOWN (alarme, exit 142) $\times3$.\\
Orbites, sib / one / dis00 / dis10 & Kissat natif & 8 h & UNKNOWN $\times4$.\\
Recuit bit-flip, graines 1,2,3 & --- & 2643 s / 85 s / 1099 s (instant de la dernière amélioration ; runs arrêtés sans bilan JSON) & meilleures énergies 1550 / 1692 / 1568 (composantes $(R,D,I,P)$ : $(14,16,270,1250)$ pour la graine 1).\\
Recuit bit-flip, graine 5 / 99 & --- & 600 s / 36 s & 1656 / 1832.\\
Recuit structuré, graines 1,2,3 / 9 & --- & 600 s / 30 s & 1966 / 2020 / 1890 / 2226 (moins bon que le bit-flip).\\
DFS \code{joint\_dfs}, \code{sout\_dfs} & --- & --- & bloquent aux lignes 2--6 (mémoire) ; pas de trace archivée.\\
\bottomrule
\end{longtable}
Les violations dominantes dans les recuits sont celles de (P) ($\approx80\,\%$ de l'énergie), puis (I) ; (R) et (D) sont presque satisfaites. Aucun de ces calculs n'apporte d'évidence dans un sens ou dans l'autre.

\section{Audit du recensement d'ordre 7 (arXiv:2608.19410v1) \st{C}}\label{sec:audit}
Méthode : transcription exacte (sympy) des 208 fréquences $z_i$, toutes affines en $(n_3,z_{11})$ ; test de l'identité $\sum_iz_i=\binom n7$ sur le lieu $n=k(k-2)/2+k+1$ ; appariement des valeurs à $n_3=705$ aux 208 classes de la famille affine du dépôt public (vague 23, système de rang 207) ; reconstruction exacte de cette famille à $n_3$ symbolique.
Constats : telles qu'imprimées, les formules laissent un déficit constant $\sum z_i-\binom n7=-\tfrac1{168}nk(k-2)(k-4)(19k^4-456k^3+4484k^2-21752k+42008)$, soit $-94\,034\,160$ à $(99,14)$ ; exactement quatre formules ne correspondent à aucune classe et quatre classes ne sont atteintes par aucune formule. Corrections : $z_{21}$ (supprimer le facteur $(k-3)$), $z_{80},z_{84}$ ($(2k-7)\to(2k-17)$), $z_{176}$ ($\tfrac1{84}\to\tfrac18$). Après correction l'identité est polynomiale en $k$ et les 208 formules coïncident avec les 208 classes sur les trois coefficients. Le recensement corrigé n'impose que $0\le n_3\le4158$, $2n_3\le z_{11}\le4n_3$, $4\mid z_{11}$ : rien de nouveau sur $n_3$.
\textbf{Reproductibilité} : les scripts exigent le source \TeX{} d'arXiv et un clone du dépôt public, absents d'ici ; les fichiers JSON produits suffisent à revérifier les identités finales. \emph{Révision 2 : fait.} À partir de \code{spec99.json} (formules imprimées) et des quatre corrections, la somme corrigée vaut exactement $\binom{99}{7}$ (coefficients de $n_3$ et $z_{11}$ nuls), le multi-ensemble des 208 triplets $(a,b,4c)$ corrigés est \emph{égal} à celui de \code{repo\_family\_affine.json}, et l'ensemble admissible recalculé ($0\le n_3\le4158$, $z_{11}\in[\lceil 2n_3\rceil_4,4n_3]$, pas 4) coïncide exactement avec \code{feasible\_n3\_z11.json}. L'auteur n'a pas été contacté.

\section{Anomalies, lacunes et corrections à apporter au dépôt}\label{sec:anomalies}
\begin{enumerate}[nosep]
\item \st{?} \textbf{Trois logs non documentés.} \code{results/batch\_kissat\_one10.log} contient \code{s UNSATISFIABLE / EXIT 20} ; \code{batch\_kissat\_one00.log} et \code{batch\_kissat\_dis.log} contiennent \code{Killed / EXIT 137}. Ils ont été ajoutés au \emph{premier} commit (19:00:45), dix minutes après le lancement de la batterie (18:50), ne figurent pas dans le bilan de la nuit, et aucun script du dépôt ne génère les CNF « one00/one10/dis00/dis10 ». Reconstruction la plus plausible : sous-découpage sur l'adjacence de $u_0$ à la sœur $\{(1,0),(2,1)\}$ ; dans le cas one ce sous-cas est vide par (I) (bloc 1 saturé par $D(U_0)$), d'où un UNSAT en quelques minutes, sans contenu ; les deux autres ont été tués pour relancer « one » entier et « dis00/dis10 ». \textbf{Ce log ne doit pas être cité comme un résultat} tant que la CNF correspondante n'est pas regénérée et documentée.
\emph{Révision 2 :} l'argument « vide par (I) » est confirmé indépendamment ($u_0=\{(1,0),(2,0)\}$, $D(u_0)=\{\{(1,0),(3,0)\},\{(1,1),(3,1)\}\}$ occupe les deux places du bloc 1 ; les deux membres de l'orbite sœur contiennent un sommet du bloc 1).
Les empreintes SHA-256 des CNF regénérées (\code{k14\_sib/one/dis/none}) sont dans \code{LAB\_INSTRUCTIONS.md} pour le test sur la machine de calcul.
\item \st{?} \code{run\_batch.sh} pointe vers un chemin absolu d'un autre répertoire de travail et lance 6 instances, alors que le bilan en résume 7 ; le mode \code{-q} a supprimé toutes les statistiques (le bilan le note).
\item \st{?} Sorties non archivées : \code{rowfix\_seed*.json}, \code{result\_cpsat\_none.json}, \code{rows\_U0.npy} (nombre de lignes candidates de $U_0$), traces de \code{joint\_dfs}/\code{sout\_dfs}, résultats \code{probe99} à $k=14$. Les énoncés correspondants du texte reposent sur la mémoire de session.
\item \st{?} Rédaction : \code{README} et \code{NOTE\_involution\_case.md} présentaient la condition $\F_2$ comme dérivée de la « structure 2-adique » et comme « information véritablement supplémentaire » (remarque~\ref{rem:statut}). \emph{Corrigé en révision 2.}
\item \st{?} L'audit d'ordre 7 dépend de deux entrées externes non incluses (§\ref{sec:audit}).
\item \st{?} \emph{(Révision 2)} Les contrôles d'origine ($k=4$, $k=6$) sont trop faibles pour détecter une erreur de modèle ; remplacés par le contrôle $\srg(243,22,1,2)$ (\code{control\_bvls.py}).
\item \st{?} \emph{(Révision 2)} Le dépôt est extrait avec \code{core.autocrlf=true} : les fichiers de travail sous Windows sont en CRLF. Toute empreinte SHA-256 (CNF, sommes de contrôle) doit être calculée sur des fichiers LF (générer les CNF sous Docker/WSL, ou comparer après \code{tr -d '\textbackslash r'}).
\item \st{?} \emph{(Révision 2)} Les rapports anglais et français et la note de travail présentaient la valeur 20 comme « information véritablement supplémentaire » et la batterie de nuit comme « en cours » ; corrigé.
\item Points vérifiés sans réserve : les lemmes de rigidité, le théorème d'équivalence (relu terme à terme contre \code{orbit\_sat.py}, et équivalence des formes $(a,b)$ et $(Q,S)$ de \code{joint\_dfs.py}), les spectres, le théorème $\F_2$, l'ordre $15\,360$ du stabilisateur, l'exhaustivité du découpage, la validité du sous-découpage dis, les corrections de l'audit (identités finales, revérifiées hors ligne), la taille des CNF à $m=7$ (regénérées).
\end{enumerate}

\section{Ce qu'un chercheur devrait vérifier (par ordre d'effort croissant)}\label{sec:verif}
\begin{enumerate}[nosep]
\item \textbf{(10 min)} Lemmes~\ref{lem:flip}, \ref{lem:noself}, proposition~\ref{prop:D}. Élémentaires ; toute la suite en dépend.
\item \textbf{(30 min)} Théorème~\ref{thm:modele} : l'énumération des sept types de paires, et le fait que « $A[u,\sigma u]=0$ » est bien une hypothèse justifiée et non une variable.
\item \textbf{(20 min)} Théorème~\ref{thm:F2} et sa remarque de statut ; propositions~\ref{prop:specQout}, \ref{prop:F3} et remarque~\ref{rem:c4} (\st{N} ; relues par un second audit automatique, jamais par un humain).
\item \textbf{(30 min)} Le découpage en cas et le sous-découpage dis ; la reconstruction du « one10 » (§\ref{sec:anomalies}).
\item \textbf{(1 h)} Rejouer \code{orbit\_sat.py --m 2 --solve} et \code{--m 3 --solve}, \code{control\_bvls.py} (léger), puis regénérer les trois CNF et vérifier qu'elles ont bien 666\,183 variables et 1\,510\,448 clauses (fait le 22/09).
\item \textbf{(accès institutionnel)} Vérification~\ref{ver:biblio} sur Behbahani--Lam 2011 et Crnković--Maksimović 2020.
\item \textbf{(1 h, avec les figures de l'article)} Confirmer les quatre coquilles de l'audit d'ordre 7 sur les dessins des classes 21, 80, 84, 176, et prévenir l'auteur.
\end{enumerate}

\section{Ce que nous aimerions voir démontré ou calculé}\label{sec:demontrer}
Classé par impact sur la branche involution. Pour chaque question : pourquoi, et ce que cela débloque.

\begin{question}[décider la branche par un générateur ordonné]\label{q:gen}
Construire les tables $(a,b)$ orbite par orbite (ou la matrice $\Qout$ puis les signes), avec : entrelacement des sous-matrices principales de $\Qout$ contre $\{12,(-2)^6,3^{20},(-4)^{15}\}$ et de $\Sout$ contre $\{0^7,3^{20},(-4)^{15}\}$ ; les contraintes linéaires $N_b^{\mathsf T}\Qout=4J-2N_b^{\mathsf T}$ et $C^{\mathsf T}\Sout=0$ ; le couplage $D$ ; les filtres de rang $\F_2$ ($\le20$) et $\F_3$ ($\le15$) sur les sous-matrices ; rejet d'isomorphes sous $S_2\wr S_7$ (nauty). \emph{Pourquoi} : c'est la seule voie identifiée dont l'issue est certifiable dans les deux sens (« $|\Aut|\in\{1,3\}$ » ou un graphe de Conway). \emph{Débloque} : la moitié restante du programme de symétrie.
\end{question}

\begin{question}[exclure un cas à la main]\label{q:cas}
Le cas \emph{sœur} ($D(U_0)=(1,2,1)$) est le plus contraint : $u_0$ n'a alors aucun autre voisin sur les blocs 1 et 2, et ses dix autres voisins forment, sur les blocs $3..7$, un multigraphe 4-régulier à 10 arêtes couplé en $5K_2$ dans $N(u_0)$. Existe-t-il un argument de comptage qui exclue ce cas, ou qui force la structure de $D$ relativement à la fibration sur $T(7)$ (par exemple : $D$ ne couple jamais deux sœurs, ou couple toujours des orbites de blocs disjoints) ? \emph{Débloque} : réduit d'un tiers les instances, et surtout indique quelle structure cachée porter dans le générateur.
\end{question}

\noindent\emph{Priorité (révision 2).} La question~\ref{q:orbit} est la prochaine tâche de calcul : c'est la moins chère des voies qui peuvent trancher, et elle est de toute façon la première étape de la question~\ref{q:gen}. Les questions~\ref{q:parite} et \ref{q:code} passent en dernier.
\begin{question}[matrices d'orbites à la Behbahani--Lam pour $p=2$]\label{q:orbit}
Le problème « $\Qout\in\{0,1,2\}^{42\times42}$, une entrée 2 et dix entrées 1 par ligne, sommes par bloc, $\Qout^2+\Qout=12I+4J-\Sh$ » est un problème de matrice d'orbites ordinaire, avec 8 lignes forcées sur 50. Les outils existants (Behbahani--Lam, Crnković--Maksimović) peuvent-ils l'énumérer ? Combien y a-t-il de telles matrices à équivalence près ? \emph{Débloque} : si zéro, la branche est close sans signes ; si peu, le relèvement (signes puis $99$ sommets) est facile.
\end{question}

\begin{question}[congruences par comptage d'orbites contre les recensements]\label{q:parite}
Avec les figures des 208 classes d'ordre 7 (et des classes d'ordre 4 à 6), calculer pour chaque classe le nombre de copies $\sigma$-invariantes (un point fixe $x$ plus trois orbites parmi 49, soit 18\,424 ensembles) en fonction de la structure forcée, et comparer sa parité à celle de la formule affine corrigée. Toute classe dont la formule a une parité constante (par exemple $z_{12}=3465$, le moulin à trois pales, dont les 35 copies invariantes sont bien en nombre impair) fournit un test ; une contradiction tuerait la branche. \emph{Réserve (révision 2)} : sauf pour les classes dont les copies invariantes sont entièrement dans la partie forcée (comme $z_{12}$, qui est donc un simple test de cohérence), le nombre de copies invariantes dépend de l'inconnue $H$ ; on obtient des contraintes de parité sur $H$, pas une contradiction immédiate. À l'ordre 4 tout est constant (remarque~\ref{rem:c4}). Basse priorité. \emph{Débloque} : au mieux des clauses de parité globales pour le SAT.
\end{question}

\begin{question}[le contenu $7$-adique]\label{q:7}
Sur $\Z_7$, $(S-3I)^2=-7(S-3I)$ ; on sait seulement $7\le\rg_{\F_7}(S-3I)\le22$ et $\rg_{\F_7}(A-3I)\ge26$ (dépôt public). Déterminer la forme normale de Smith $7$-adique de $S-3I$ et de $Q-3I$, et la relier au $7$-rang de $A-3I$ (Brouwer--van Eijl) via $\Z_7^{99}=\Lambda_e\oplus\Lambda_o$ (2 est inversible). \emph{Débloque} : le seul invariant $p$-adique non épuisé par la proposition~\ref{prop:F3}.
\end{question}

\begin{question}[code LCD]\label{q:code}
Classer ou borner les codes LCD $[42,20]$ contenus dans $(\operatorname{col}N_b)^\perp$ dont le projecteur $K$ vérifie : $K+N_bN_b^{\mathsf T}$ est $0/1$, à diagonale nulle, 10-régulière (remarque de §\ref{sec:demontre}). Les identités de MacWilliams excluent-elles cette configuration ? \emph{Débloque} : une obstruction purement mod 2, ou une énumération par sous-espaces plutôt que par graphes.
\end{question}

\begin{question}[cube-and-conquer]\label{q:cube}
Compter les lignes valides de $U_0$ à équivalence près sous $\operatorname{Stab}(U_0)$ (ordre 15\,360 ; \code{rows.py} les énumère sans quotienter) ; le dépôt estime $10^4$--$10^5$ lignes brutes. Avec BreakID pour casser la symétrie résiduelle et un journal DRAT par cube, une campagne certifiée est-elle réaliste (jours-cœurs) ? \emph{Débloque} : une preuve UNSAT vérifiable indépendamment (drat-trim), à défaut d'un générateur.
\end{question}

\begin{question}[une identité trace--rang plus fine]\label{q:ishida}
La méthode d'Ishida exclut les involutions du graphe de Moore manquant mais redonne seulement $a_1(t)=14$ ici. Existe-t-il une version « forme quadratique » ou « réseau libre $\Z_2[C_2]$ » qui utilise le couplage $D$ (par exemple les 21 cycles $C_4$ invariants) ? \emph{Débloque} : une exclusion sans ordinateur.
\end{question}

\begin{question}[antériorité]\label{q:ant}
Quelqu'un a-t-il déjà calculé le cas $p=2$ de $(99,14,1,2)$, même partiellement (matrices d'orbites, ou dans l'analyse de $\Z_6$) ? \emph{Débloque} : éviter de refaire, ou s'appuyer sur un code existant.
\end{question}

\section{Ressources nécessaires}\label{sec:besoins}
\begin{itemize}[nosep]
\item \textbf{Documents} : Behbahani--Lam 2011 (Discrete Math.\ 311) ; Crnković--Maksimović 2020 (Contrib.\ Discrete Math.\ 15) ; les figures de Reimbayev 2026 ; éventuellement Wilbrink 1984 (ordre 11) et la traduction anglaise de Makhnev--Minakova.
\item \textbf{Outils} : nauty/Traces sous Linux (pynauty n'a pas pu être installé sous Windows) ; BreakID ou shatter ; drat-trim ou cake\_lpr ; CaDiCaL/Kissat (déjà dans l'image \code{satbox}), éventuellement Gimsatul.
\item \textbf{Calcul} : une machine Linux 32--64 cœurs pendant plusieurs jours, pour Q\ref{q:gen} ou Q\ref{q:cube}. Les CNF se regénèrent en 3 s. \emph{(Révision 2 : la machine de calcul disponible a 8 cœurs / 16 fils et 32 Go ; elle suffit pour Q\ref{q:orbit}, le contrôle \code{control\_bvls.py --cnf} et des essais de Q\ref{q:gen}.)}
\item \textbf{Humain} : la liste de §\ref{sec:verif}, en priorité les points 1--4 (deux heures au total).
\end{itemize}

\appendix
\section{Reproduction}
\begin{verbatim}
python -m pip install sympy numpy python-sat ortools     # Python 3.12
python code/z2/orbit_sat.py --m 2 --solve   # SAT, releve verifie srg(9,4,1,2)
python code/z2/orbit_sat.py --m 3 --solve   # UNSAT
python code/z2/orbit_sat.py --m 7 --case one --out k14_one.cnf
docker build -t satbox code/z2 && kissat --time=28800 k14_one.cnf
python code/z2/control_bvls.py              # controle fort, leger (< 2 min)
python code/z2/control_bvls.py --cnf        # machine de calcul uniquement (~20 Go)
docker build -t texbox -f report/Dockerfile.tex report && \
  docker run --rm -v "$PWD/report:/work" texbox pdflatex dossier_chercheurs_fr.tex
\end{verbatim}

\section{Bibliographie du dépôt}
\begin{itemize}[nosep]
\item A.\,A.~Makhnev, I.\,M.~Minakova, Diskretn.\ Mat.\ 16:1 (2004) 95--104, DOI 10.4213/dm144 ; trad.\ Discrete Math.\ Appl.\ 14:2 (2004) 201--210.
\item M.~Behbahani, \emph{On strongly regular graphs}, thèse, Concordia, 2009 (LAC NR63369) ; M.~Behbahani, C.~Lam, Discrete Math.\ 311 (2011) 132--144.
\item D.~Crnković, M.~Maksimović, Contrib.\ Discrete Math.\ 15 (2020) 22--41.
\item P.\,G.~Cesarz, A.\,J.~Woldar, Algebraic Combinatorics 8 (2025) 379--398.
\item D.~Thakkar, S.~Severini, arXiv:2608.11211v2 (2026). Y.~Ishida, arXiv:2606.29183v2 (2026). R.~Reimbayev, arXiv:2608.19410v1 (2026).
\item YesterdaysLemon/conway-99-research (GitHub, 2026).
\item A.\,A.~Makhnev, \emph{Graphs and automorphisms}, notes de l'école d'été de Bielefeld (\code{math.uni-bielefeld.de/\textasciitilde baumeist/sommerschule/makhnev.pdf}).
\item A.~Keramatipour, \emph{Approaching the Conway-99 problem using SAT solvers}, arXiv:2604.23037 (2026).
\item E.\,R.~Berlekamp, J.\,H.~van Lint, J.\,J.~Seidel, A strongly regular graph derived from the perfect ternary Golay code, dans \emph{A Survey of Combinatorial Theory} (1973) 25--30 (contrôle $\srg(243,22,1,2)$).
\item A.\,E.~Brouwer, C.\,A.~van Eijl, \emph{On the p-rank of the adjacency matrices of strongly regular graphs}, J.\ Algebraic Combin.\ 1 (1992) 329--346 (pour Q\ref{q:7}).
\end{itemize}
\end{document}
